\subsubsection{B4.1 Zwei zusätzliche Terme am unveränderten
Tensor}\label{b4.1-zwei-zusuxe4tzliche-terme-am-unveruxe4nderten-tensor}

Der Bosonträger (10,6) besitzt eine symmetrische reelle invariante
Paarung eta. In der nativen Gewichtsbasis koppelt sie entgegengesetzte
Spin- Vektorgewichte und komplementäre Farbpaare. Es gilt eta Quadrat
gleich eins. Damit sind bereits

\[\adjustbox{max width=\linewidth}{$\displaystyle 
B_+=\tfrac12\sum_{AB}b_A^\dagger\eta_{AB}b_B^\dagger,
\qquad
R_+=\sum_{AB}b_A^\dagger\eta_{AB}P_B^\dagger
$}\]

G-invariant. Bplus erzeugt zwei Bosonen, Rplus einen Boson und zwei
Fermionen. Beide ändern N um vier; beide erhalten N modulo vier. Die
vollständige Lie-Invarianz wurde an allen 45 plus 15 Generatoren
überprüft, nicht nur an den Cartanladungen. Rplus wirkt auf dem leeren
Zustand nicht trivial: Das Normquadrat beträgt exakt 480. Der
quadratische Bosonpaarterm ist sogar einfacher als der vorgeschlagene
zusätzliche kubische Kanal.

Auf dem unveränderten Träger lautet die allgemeine G-invariante,
fermionparitätsgerade, normalgeordnete Hamiltonfamilie vom Grad
höchstens drei

\[\adjustbox{max width=\linewidth}{$\displaystyle 
H_{\rm ext}=c+\varepsilon N_f+\Delta N_b
+\kappa B_++\bar\kappa B_-
+g\sum_A b_A^\dagger P_A+\bar g\sum_A P_A^\dagger b_A
+\lambda R_++\bar\lambda R_-.
$}\]

Die Vollständigkeit gilt nur innerhalb dieser ausdrücklichen
Polynomialklasse: End(F) und End(B) haben jeweils einen Skalar, Sym hoch
zwei von B einen Skalar, und B tritt in Lambda hoch zwei von F genau
einmal auf. Konkret zerfällt letzterer Raum in (10,6), (126,6) und
(120,10). Die SU(4)-Zentrumsregel schließt die verbleibenden
bosonisch-kubischen und gemischten Besetzungs-Boson-Terme aus. Höhere
Polynome, räumliche Kopien und neue Feldträger sind damit nicht
klassifiziert.

Für das native Modell hat das lineare System aller diagonalen
Modenphasen exakt neun unabhängige kontinuierliche Erhaltungen: die acht
inneren Cartans und N. Bei zusätzlichem nichtverschwindendem Rplus-Kanal
oder Bosonpaar-Kanal bleiben exakt acht. Die Ränge sind 115
beziehungsweise 116 auf 124 Moden; explizite ganzzahlige Kerne und
modulare Rangzertifikate beweisen sie. Dies ist keine Behauptung, die
volle nichtabelsche Gruppe schrumpfe auf acht Dimensionen. Sie bleibt
erhalten. Außerdem ist N modulo vier bereits die Wirkung des
SU(4)-Zentrums und keine neu hinzugewonnene unabhängige Symmetrie.

Für reelle Kopplungen und epsilon gleich null enthält die
Fermiongleichung nun zwei verschiedene Beiträge:

\[\adjustbox{max width=\linewidth}{$\displaystyle 
[H_{\rm ext},f_r]=-gD_r-\lambda\widetilde\chi_r^\dagger,
\qquad D_r\sim bf^\dagger,\quad
\widetilde\chi_r^\dagger\sim b^\dagger\eta f^\dagger.
$}\]

Sie haben N-Ladungen -1 und +3. Eine Dynamik kann diese Kanäle bei
bloßer Z4-Erhaltung verbinden, ohne die Operatoren oder ihre
Referenzzustände fälschlich gleichzusetzen. \textbf{Zulässig ist aber
noch nicht aus dem Compiler abgeleitet.} Die Grundzustands- und
Polschranken des unveränderten Modells gelten für diese Erweiterung
nicht ungeprüft weiter.

Auch ein mu-N-Term bleibt G- und Z4-invariant. Bereits der
Zwei-Zustandsblock aus leerem Zustand und einer normierten
Quartetterzeugung im kleinen Modell zeigt

\[\adjustbox{max width=\linewidth}{$\displaystyle 
H+\mu N=\begin{pmatrix}0&\lambda\\\lambda&\Delta+4\mu\end{pmatrix}.
$}\]

Seine niedrigere Eigenenergie hängt von mu ab. Der Übergang zu Z4 allein
wählt also weder ein chemisches Potential noch einen eindeutigen
physikalischen Grundzustand aus. Für Delta größer als Betrag kappa ist
die bosonische Quadratik positiv; die linearen Bosonkopplungen an
beschränkte Fermionoperatoren lassen sich relativ dazu abschätzen. Das
gibt eine untere Energieschranke, aber noch keinen neuen eindeutigen
Grundzustandssatz.

\subsubsection{B4.2 Eine weitere Vereinfachung: neue Dynamik oder andere
Teilchenvariablen?}\label{b4.2-eine-weitere-vereinfachung-neue-dynamik-oder-andere-teilchenvariablen}

Manche scheinbaren Erweiterungen beschreiben nur andere Bosonvariablen.
Auf der reellen, phasengleichen Teilfamilie setze man

\[\adjustbox{max width=\linewidth}{$\displaystyle 
b=c\cosh r+\eta c^\dagger\sinh r.
$}\]

Dies ist eine G-verträgliche kanonische Bogoliubov-Transformation. Sie
ergibt

\[\adjustbox{max width=\linewidth}{$\displaystyle 
\Delta'=\Delta\cosh2r+\kappa\sinh2r,\quad
\kappa'=\Delta\sinh2r+\kappa\cosh2r,
$}\] \[\adjustbox{max width=\linewidth}{$\displaystyle 
g'=g\cosh r+\lambda\sinh r,\quad
\lambda'=g\sinh r+\lambda\cosh r.
$}\]

Die additive Konstante ist 30 mal (Delta prime minus Delta). Bei g
Quadrat größer lambda Quadrat verschwinden kappa prime und lambda prime
gemeinsam genau dann, wenn

\[\adjustbox{max width=\linewidth}{$\displaystyle 
\mathcal I=\kappa(g^2+\lambda^2)-2\Delta g\lambda=0.
$}\]

Beispiel: Delta=1, kappa=4/5, g=2, lambda=1 und tanh r=-1/2 ergeben
Delta prime=3/5, g prime=Wurzel 3, kappa prime=lambda prime=0. Obwohl
die alte Zahl N nicht erhalten ist, gibt es dann die verborgene
kontinuierliche Erhaltung Nf plus zweimal Nc. Daher wäre auch die
Folgerung „N ist gebrochen, also ist das Fundament nur Z4`` ohne Prüfung
zu schnell. Für I ungleich null ist nur diese uniforme kanonische
Reduktion ausgeschlossen, nicht jede denkbare nichtlineare verborgene
Symmetrie.

Der sehr einfache reelle Quadraturansatz lambda=g bei kappa=0 ist
dagegen nicht durch eine endliche solche Transformation auf den nativen
Ansatz reduzierbar. Er ist ein prüfbarer Kandidat, kein hergeleitetes
Naturgesetz. Ein weiterer Herkunftstest ist besonders schlicht: Bei
freiem H0=epsilon Nf+Delta Nb haben die beiden kubischen Terme die
Frequenzen Delta minus 2 epsilon und Delta plus 2 epsilon. Am
verwendeten epsilon=0 gibt es keine schnelle gegen langsame Frequenz,
die das Weglassen des zweiten Kanals durch eine Rotating-Wave-Näherung
rechtfertigen würde. Ein exaktes U(1)-Prinzip könnte das dennoch
rechtfertigen; dieses Prinzip wäre dann auszuweisen und abzuleiten.

\subsection{B5. Holografie und Schatten: die konkrete native
Verbindung}\label{b5.-holografie-und-schatten-die-konkrete-native-verbindung}

\subsubsection{B5.1 Ein Kodierer ist bereits
vorhanden}\label{b5.1-ein-kodierer-ist-bereits-vorhanden}

Der vorhandene, frisch geprüfte Tensor erfüllt WW dagger gleich 8 I60.
Daher ist

\[\adjustbox{max width=\linewidth}{$\displaystyle 
V=W^\dagger/\sqrt8:\ \mathbb C^{60}\longrightarrow\Lambda^2\mathbb C^{64},
\qquad V^\dagger V=I_{60},\quad
\Pi=VV^\dagger=W^\dagger W/8
$}\]

eine exakte Isometrie mit einem Rang-60-Projektor Pi. Ein
60-dimensionaler logischer Zustand kann verlustfrei als kohärente
Überlagerung von Fermionpaaren dargestellt werden. Die andere Richtung
sieht aber nicht den ganzen 2016-dimensionalen Paarraum: Ihr Kern hat
Dimension 1956. Diese „dunklen`` Richtungen sind nicht deshalb generell
unphysikalisch; sie sind für genau diese Abbildung unsichtbar.

Die beiden Ansichten sind durch W also bereits verbunden. Die native
N=2-Dynamik im hellen Paar-plus-Boson-Bereich ist sogar exakt

\[\adjustbox{max width=\linewidth}{$\displaystyle 
H_{N=2,\mathrm{hell}}=
\begin{pmatrix}0&\sqrt8 g\\\sqrt8 g&\Delta\end{pmatrix}\otimes I_{60}.
$}\]

Sie wandelt die Kodierungsform um; am Prüfpunkt erreicht die
entsprechende Übergangswahrscheinlichkeit höchstens 2/27. Der nur in die
Paarseite eingebettete Code ist bei g ungleich null kein invarianter
H-Unterraum. Dieser kleine gemeinsame dynamische Träger ist ein
positiver Befund. Er ist N=2, nicht der native N=64-Grundzustand und
nicht dessen geladener Pol. Insbesondere ist er noch kein Transport
durch physikalischen Raum.

\subsubsection{B5.2 Warum das noch kein holografischer
Fehlerkorrekturcode
ist}\label{b5.2-warum-das-noch-kein-holografischer-fehlerkorrekturcode-ist}

Für jeden der 64 Fermionmoden wurde die logische Kompression seiner
Besetzungsmessung berechnet:

\[\adjustbox{max width=\linewidth}{$\displaystyle 
V^\dagger n_rV\quad\hbox{hat Spektrum}\quad
0\ (45\text{-fach}),\quad 1/8\ (15\text{-fach}).
$}\]

Das Resultat ist nicht skalar. Die Umgebung kann daher schon durch
Ablesen einer einzigen verlorenen Mode etwas über den logischen Zustand
erfahren. Die notwendige Fehlerkorrekturbedingung für die vollständige
60er-Codemenge ist verletzt: Der Code korrigiert nicht die beliebige
Löschung dieser einen Mode. Fermionparität ändert dieses
Besetzungsargument nicht, denn n\_r ist gerade. Kleinere geeignete
Untercodes oder ein aus nativen Operationen erzeugtes größeres
Kodierungsnetz sind dadurch nicht ausgeschlossen, aber noch nicht
konstruiert. Bloßes Aneinanderhängen desselben Tensors garantiert keinen
Code.

In etablierten holografischen Spielzeugmodellen sind isometrische
Kodierung, Rekonstruktion logischer Operationen aus verschiedenen
Randteilen und Fehlerkorrektur genau kontrollierte Eigenschaften. Das
ist ein hilfreiches Prüfschema, keine direkte Identifikation mit TFPT.
Siehe \href{https://arxiv.org/abs/1503.06237}{Pastawski, Yoshida, Harlow
und Preskill}. Die Rekonstruktion aus einem Teilrand ist auch in der
AdS/CFT-Arbeit von \href{https://arxiv.org/abs/1601.05416}{Dong, Harlow
und Wall} an konkrete Quanteninformations- und geometrische
Voraussetzungen gebunden. Diese Voraussetzungen wurden für TFPT nicht
nachgewiesen.

\subsubsection{B5.3 Was mehrere Schatten gemeinsam eindeutig machen
können}\label{b5.3-was-mehrere-schatten-gemeinsam-eindeutig-machen-kuxf6nnen}

Seien Ri festgelegte Mess- oder Reduktionsabbildungen eines gemeinsamen
Zustands rho. Für eine endliche, nicht eingeschränkte Zustandsklasse
sind alle Zustände aus den Schatten Ri(rho) genau dann unterscheidbar,
wenn auf den hermiteschen spurlosen Differenzen gilt

\[\adjustbox{max width=\linewidth}{$\displaystyle 
\bigcap_i\ker R_i=\{0\}.
$}\]

Das ist das einfache gemeinsame-Kern-Kriterium. Anschaulich: Was eine
Ansicht nicht sieht, muss eine andere sehen. Zwei Qubitansichten auf
(x,z) und (y,z) genügen gemeinsam zur Bloch-Rekonstruktion, jede für
sich nicht. Zweimal dieselbe Ansicht auf (x,z) genügt weiterhin nicht.
Das gilt ebenso für zwei Theorieberichte, deren Übereinstimmung bereits
aus denselben eingesetzten Annahmen stammt: Sie liefern keine zweite
unabhängige Messung.

Drei exakt nachgerechnete Grenzen sind wesentlich:

\begin{enumerate}
\def\labelenumi{\arabic{enumi}.}
\tightlist
\item
  Die orthogonalen Zustände (000 plus 111)/Wurzel 2 und (000 minus
  111)/Wurzel 2 haben dieselben reduzierten Zustände auf sämtlichen
  echten Teilmengen der drei Qubits. Eine gemeinsame phasensensitive
  XXX-Messung unterscheidet sie mit Erwartungswert +1 oder -1. Auch sehr
  viele lokale Schatten können also eine globale Phase übersehen.
\item
  Paarweise passende Überschneidungen garantieren nicht einmal einen
  gemeinsamen Ursprung. Drei binäre Variablen können nicht paarweise
  ausnahmslos verschieden sein. Jede einzelne perfekte Antikorrelation
  besitzt aber dieselben gleichverteilten Einzelmarginalen. Gemeint ist
  hier ein gemeinsames klassisches Tripel, nicht eine Behauptung gegen
  kontextuelle Quantenexperimente.
\item
  Selbst die vollständige Antwort einer Referenz bestimmt keinen völlig
  entkoppelten dunklen Zusatzsektor. H und H direkt plus K besitzen vom
  Startvektor (F,0) aus dieselben Energiemomente und dieselbe
  Zeitantwort. Ein exaktes kleines Matrixbeispiel prüft diesen
  Sachverhalt.
\end{enumerate}

Das erreichbare Ziel ist daher zunächst der \textbf{kleinste gemeinsame
beobachtbare Prozess}, nicht ein aus endlichen Schatten bewiesener
einzigartiger ontologischer Innenraum. Zwei Historien gelten operational
als gleich, wenn alle verfügbaren künftigen Eingriffs- und Auslesefolgen
dieselben Wahrscheinlichkeiten liefern. Das ist eine eindeutige
Unterscheidungsregel für gegebenes Verhalten; sie garantiert weder eine
eindeutige Hilbertraumdarstellung noch von selbst eine einfache
Geometrie.

\subsubsection{B5.4 Die drei Wände müssen unterschiedlich geprüft
werden}\label{b5.4-die-drei-wuxe4nde-muxfcssen-unterschiedlich-gepruxfcft-werden}

{\def\LTcaptype{none} % do not increment counter
\begin{longtable}[]{@{}
  >{\raggedright\arraybackslash}p{(\linewidth - 4\tabcolsep) * \real{0.3333}}
  >{\raggedright\arraybackslash}p{(\linewidth - 4\tabcolsep) * \real{0.3333}}
  >{\raggedright\arraybackslash}p{(\linewidth - 4\tabcolsep) * \real{0.3333}}@{}}
\toprule\noalign{}
\begin{minipage}[b]{\linewidth}\raggedright
Gemeinte Wand
\end{minipage} & \begin{minipage}[b]{\linewidth}\raggedright
Was sie gegenwärtig bedeutet
\end{minipage} & \begin{minipage}[b]{\linewidth}\raggedright
Welcher Anschluss fehlt
\end{minipage} \\
\midrule\noalign{}
\endhead
\bottomrule\noalign{}
\endlastfoot
Universalraum & Grenze zwischen gemeinsamem Operatorobjekt und seinen
reduzierten Ansichten & Ein gemeinsamer verfügbarer Operationssatz,
Zustand und Rekonstruktionsvertrag \\
TFPT & Grenze zwischen interner Compilerstruktur und physikalischer
Realisierung & Herleitung der Ausführung, relativer Parameter und
operationaler Teilung \\
Beobachtete Realität & Endlicher Zugang über Messungen, Präparationen
und Korrelationen & Ein quantitatives, nicht nachträglich angepasstes
Feld- und Messwörterbuch \\
\end{longtable}
}

Diese Erkenntnisgrenzen sind nicht schon drei bewiesene physikalische
Holografieschirme. Ein kosmologischer Horizont, ein Code-Rand und eine
offene mathematische Beweispflicht sind verschiedene Dinge. Ebenso sind
Universalraum und TFPT Kandidatenbeschreibungen der Realität, nicht
bereits drei experimentell etablierte wechselwirkende Welten.

Eine stärkere holografische Deutung müsste eine einzige Kodierung mit
Zuständen, Operationen und Zeitentwicklung verbinden. Auf einem
invarianten Codesektor wäre etwa V dagger V=I und Hrand V=V Hin eine
passende Intertwining-Bedingung. Bei nichtinvarianten reduzierten
Ansichten muss die entstehende Erinnerung mitgeführt werden; autonome
reduzierte Dynamik darf nicht einfach vorausgesetzt werden. Zusätzlich
braucht es eine Bedeutung von Randteilen, kontrollierte Rekonstruktion
und eine skalierende Geometrie.

\subsection{B6. Die einfachste neue Arbeitsrichtung und ihre
Abnahmekriterien}\label{b6.-die-einfachste-neue-arbeitsrichtung-und-ihre-abnahmekriterien}

Der plausible übersehene Punkt ist kein weiterer großer Tensor:
\textbf{Wir brauchen die gemeinsame Regel, die festlegt, welche
Operationen, Teilchenvariablen und Ansichten tatsächlich dieselbe
Ausführung beschreiben.} Noethers Symmetrieprüfung sagt, was erhalten
bleibt. Die Schattenprüfung sagt, was überhaupt unterscheidbar ist. Der
Eingriffstest sagt, was etwas anderes beeinflussen kann. Zusammen sind
diese drei kleinen Fragen strenger und informativer als weitere
isolierte Zahlenübereinstimmungen.

\begin{enumerate}
\def\labelenumi{\arabic{enumi}.}
\tightlist
\item
  \textbf{Den minimalen Ursprungsvertrag auswählen.} Den tatsächlichen
  Compiler auf den nativen kubischen Term, den zweiten kubischen Kanal
  und den quadratischen Paarterm prüfen. Die vollständige kleine Familie
  und der Bogoliubov-Test verhindern, dass derselbe Prozess mehrfach
  gezählt wird. Erfolg ist eine Ableitung oder ein präziser Ausschluss
  samt relativen Parametern, nicht nur die Feststellung erlaubter
  Symmetrie.
\item
  \textbf{Die Schattenkarte am vorhandenen W aufbauen.} Für denselben
  Vertrag Präparationen, Aufzeichnungen und kontrollierte
  Mehrzeitantworten in den verschiedenen Ansichten auflisten. Gemeinsam
  unsichtbare Richtungen und unverträgliche Überlappungen berechnen. Aus
  einem Teil der Daten eine Aussage bestimmen, die eine andere Ansicht
  ohne Nachjustierung testet. Der jetzt bewiesene Rang-60-Code und sein
  Einmoden-Leck sind Ausgangsdaten, kein Anlass, Raum oder perfekte
  Fehlerkorrektur bereits anzunehmen.
\item
  \textbf{Daraus einen gemeinsamen kausalen Teilträger gewinnen.} Eine
  verfügbare Operation muss die unbedingte Statistik eines anderen
  operational bestimmten Teils verändern. Derselbe Träger muss das
  relativistische Feldwörterbuch tragen. Erst dann sind Skalierung,
  chirales Maß und ein Spin-2-Sektor belastbare nächste Rechnungen.
\end{enumerate}

Die lokale Lanczos-Fortsetzung bleibt als kontrolliertes
Diagnoseinstrument nützlich, steht aber nicht mehr an erster Stelle der
fundamentalen Suche. Die drei Schritte ersetzen keine einzelnen
T1-T8-Nachweise. Auch eine erfolgreiche endliche Kodierung würde RH,
effiziente Faktorisierung oder P-versus-NP nicht automatisch lösen.

Der RH-Absatz der Eingabe enthält keine neue vollständige arithmetische
Beweiskonstruktion. Das vorhandene Forschungsregister wurde begrenzt
abgeglichen; sein Aktualitätslauf scheitert derzeit an einer nicht
verfügbaren historischen Quelle beziehungsweise Quellen-/Review-Drift.
Bestehende endliche Determinantsonden sind ausdrücklich Diagnosemodelle.
Es wurde kein neuer RH-Beweis, kein umfassender RH-Neulauf und keine
Erneuerung der entsprechenden Vertrauensmarker behauptet.

\textbf{Bilanz:} Ein gemeinsamer Kodierungsbaustein ist schon da;
außerdem sind die minimalen symmetrieverträglichen Dynamiken und zwei
wichtige Verwechslungen jetzt wesentlich schärfer bestimmt. Die
vollständige universelle physikalische Lösung bleibt offen. Der nächste
Schritt ist kleiner und konkreter geworden: dieselbe Ausführung durch
mehrere kalibrierte, dynamisch konsistente Ansichten rekonstruieren und
testen.

